\renewcommand{\theequation}{\arabic{chapter}.\arabic{equation}}

\chapter{MARCO TEÓRICO}
\label{cap:teoria}

Este capítulo presenta los conceptos que usa el método del Capítulo~\ref{cap:metodologia}, en el
orden en que los necesita. Empieza por el simulador SUMO y la interfaz TraCI, con los que se
construyen los escenarios (U1) y se miden los controles (U3). Sigue con el control semafórico, sus
dos controles de referencia, la coordinación entre cruces y las métricas con que U3 compara los
controles. Luego presenta los fundamentos del aprendizaje por refuerzo, en los que se apoya la
formulación del control (U2), y los métodos tabulares Q-learning y SARSA. Siguen los métodos con
redes neuronales, DQN y PPO, y el caso de varios agentes, que es el marco de IPPO, MAPPO y la
mensajería entre vecinos (U4). Cierra con los conceptos de evaluación del plan de pruebas.

\section{Simulación microscópica de tráfico}
\label{sec:teoria-simulacion}

\subsection{SUMO}

SUMO (\textit{Simulation of Urban MObility}) es un microsimulador de tráfico de código abierto:
modela el movimiento de cada vehículo por separado \citep{TraCI_Whitepaper}. La red vial es un
grafo dirigido cuyos nodos son intersecciones y cuyas aristas son tramos con uno o más carriles y
sus restricciones de giro. Se construye con \texttt{netconvert}, a partir de archivos de nodos y
vías o de cartografía de OpenStreetMap \citep{SUMO_Netconvert_OSM,OSM_About}. Sobre la red se
declaran los tipos de vehículo, la demanda y los programas de los semáforos. Al terminar cada
viaje, SUMO escribe un registro por vehículo (\texttt{tripinfo}), del que salen casi todas las
métricas de la Sección~\ref{sec:teoria-metricas}. El tiempo de viaje y el retraso del registro se
cuentan desde que el vehículo entra a la red. Si la vía por la que debe entrar no tiene espacio,
SUMO lo retiene fuera hasta que lo haya y anota esa espera aparte, en el mismo registro
(\texttt{departDelay}). Un vehículo que al cerrar la simulación todavía espera entrar no tiene
registro: solo lo cuenta el resumen de estadísticas de la corrida (\verb|--statistic-output|).

La conducción sigue un modelo de seguimiento vehicular con componentes aleatorias: la
imperfección del conductor, controlada por el parámetro $\sigma$, y un factor de velocidad
individual. La semilla de la simulación fija esas realizaciones: con la misma semilla y el mismo
control, dos corridas son idénticas. Por eso la semilla es la unidad de réplica de los
experimentos (Sección~\ref{sec:evaluacion}). Qué varía entre semillas en los escenarios de esta
tesis se detalla en el Capítulo~\ref{cap:metodologia}.
% TODO(reestructura): citar la documentacion de SUMO del modelo de seguimiento por defecto (Krauss, sigma,
% speedFactor); no se pudo verificar que TraCI_Whitepaper los describa. Descripcion tomada del capitulo4 original.

\subsection{TraCI}

TraCI (\textit{Traffic Control Interface}) es la interfaz cliente--servidor de SUMO: un programa
externo, en esta tesis escrito en Python, lee el valor de los objetos de la simulación y modifica su
comportamiento mientras la simulación corre \citep{TraCI_Docs,TraCI_Whitepaper}. El cliente
avanza la simulación paso a paso y, entre dos pasos, lee observaciones (vehículos detenidos por
aproximación, tiempos de espera, fase activa) y aplica acciones (fijar la fase de un semáforo o su
duración); para no multiplicar las llamadas, puede suscribirse a un conjunto de variables
\citep{TraCI_Docs}. Es el canal por el que los agentes actúan sobre el tráfico y por el que se
miden las magnitudes que no están en el registro por vehículo, como las colas en cada segundo.

\section{Control semafórico}
\label{sec:teoria-control}

\subsection{Fases, ciclo y parámetros de temporización}

Un semáforo reparte el derecho de paso entre movimientos incompatibles. Una fase es el intervalo
en que un conjunto de movimientos compatibles recibe verde a la vez; cada fase verde termina con un
intervalo de cambio en ámbar. El programa es la secuencia de fases con sus duraciones, y el ciclo,
el tiempo que tarda en recorrerla una vez. El reparto del verde (\textit{split}) es la división
del ciclo entre las fases, y el desfase (\textit{offset}), el retraso del inicio del ciclo de un
semáforo respecto de una referencia común a los semáforos coordinados \citep{HCM2000_or_2010}. En
SUMO un programa es una lista de fases, cada una con su duración y el estado de cada señal, más un
desfase inicial \citep{SUMO_TLS_Actuated}.

El verde mínimo es el tiempo que una fase debe mantenerse antes de poder cambiar, y el verde
máximo, el tiempo tras el cual el cambio se fuerza aunque siga llegando demanda
\citep{HCM2000_or_2010,SUMO_TLS_Actuated}. En esta tesis los agentes y el control actuado operan
con los mismos límites de verde y el mismo ámbar (Capítulo~\ref{cap:metodologia}).
% TODO(reestructura): verificar en el HCM 2010 la seccion que define fase, ciclo, reparto, desfase, verde minimo y
% maximo, control prefijado y actuado y progresion. El capitulo 3 del manual de la FHWA citado ("Operational and
% Safety Analysis") no los define: solo capacidad, flujo de saturacion, grado de saturacion, control delay y colas.

Durante el verde, una cola se descarga a un ritmo máximo, el flujo de saturación. La capacidad de
un carril semaforizado es ese flujo multiplicado por la fracción del ciclo en verde efectivo, y el
cociente entre demanda y capacidad es el grado de saturación
\citep{FHWA_Timing_Manual_Chapter3,HCM2000_or_2010}. Si es mayor que uno, la demanda supera lo que
el semáforo puede despachar y la cola crece.

\subsection{Control de tiempo fijo y control actuado}

El control de tiempo fijo y el actuado son los dos controles de referencia contra los que se mide
cualquier mejora. El control de tiempo fijo repite el mismo programa en todos los ciclos, con duraciones independientes del tráfico
presente. Su calidad depende del ajuste del ciclo, del reparto y, en un corredor, de los desfases
\citep{HCM2000_or_2010}. Comparar con un programa fijo sin ajustar sobrestima la
mejora de cualquier alternativa; por eso el método compara contra un fijo ajustado por barrido en
cada caso, el fijo tuneado (Capítulo~\ref{cap:metodologia}).

El control actuado conserva las fases, pero adapta cada verde a lo que miden los detectores:
prolonga la fase mientras llega un flujo continuo de vehículos y pasa a la siguiente cuando
detecta un hueco suficiente entre dos vehículos, siempre entre el verde mínimo y el máximo.
Reacciona a la demanda instantánea, sin aprender ni anticipar. SUMO lo implementa de forma
nativa, con una duración mínima (\texttt{minDur}) y una máxima (\texttt{maxDur}) por fase
\citep{SUMO_TLS_Actuated}.

\subsection{Coordinación, onda verde y \textit{spillback}}

En un corredor, los semáforos se coordinan cuando comparten el ciclo y sus desfases se eligen para
favorecer el paso de los pelotones. En la onda verde, el desfase entre dos semáforos sucesivos es
igual al tiempo de viaje entre ellos, de modo que un pelotón que sale en verde del primero
encuentra verde en los siguientes. Es la forma clásica de coordinar un corredor con control fijo
\citep{HCM2000_or_2010}; por eso es una referencia exigente para un agente que decide con
información local. La progresión se mide con la fracción de
vehículos que cruzan el corredor deteniéndose a lo sumo una vez (Sección~\ref{sec:teoria-metricas}).
% TODO(reestructura): progresion en los dos sentidos y desfases alternados de medio ciclo ("en ajedrez", fijo de
% P5 en tab:plan2-fijo): falta una fuente (capitulo de coordinacion del manual de la FHWA o el HCM).
% En el analisis de robustez (re-barrido v2.1, tab:robustez-fijo) el ajedrez queda segundo en P5 y gana 10/10 s
% con los desfases en cero; el fijo del resultado principal sigue siendo el ajedrez.

Cuando la cola de un enlace entre dos semáforos ocupa todo el espacio disponible, rebasa el enlace
y bloquea la intersección aguas arriba, que ya no puede despachar hacia él aunque tenga verde: es
el \textit{spillback} \citep{FHWA_Timing_Manual_Chapter3,HCM2000_or_2010}. La cola que un
semáforo empuja hacia su enlace de salida pertenece a las aproximaciones del vecino, no a las
suyas; penalizarla en la recompensa es una de las variantes de U4
(Capítulo~\ref{cap:metodologia}).

\section{Métricas de desempeño}
\label{sec:teoria-metricas}

Todas las métricas se calculan con las mismas salidas de SUMO para todos los controles, de modo
que la comparación no dependa de cómo decide cada uno. La Tabla~\ref{tab:metricas} las define; el
Capítulo~\ref{cap:metodologia} describe cómo se mide cada una.

\begin{table}[H]
\setstretch{1}
\centering
\footnotesize
\begin{tabular}{>{\raggedright\arraybackslash}p{2.7cm}>{\raggedright\arraybackslash}p{7.1cm}>{\raggedright\arraybackslash}p{2.4cm}}
\toprule
Métrica & Definición & Salida de SUMO \\
\midrule
Retraso (primaria) & tiempo de viaje real menos el que tomaría el mismo recorrido a la velocidad
  deseada, sin otros vehículos ni semáforos, contado desde que el vehículo entra a la red (s/veh) &
  \texttt{timeLoss} (\texttt{tripinfo}) \\
Espera para entrar & tiempo entre la salida programada y la entrada a la red, que el vehículo pasa
  fuera por falta de espacio en la vía de entrada (s/veh) & \texttt{departDelay}
  (\texttt{tripinfo}) \\
Retraso total (robustez) & retraso más espera para entrar, incluida la de los vehículos que al
  cerrar la simulación no habían entrado (s/veh; Ecuación~\ref{eq:retraso-total}) &
  \texttt{timeLoss} y \texttt{departDelay} (\texttt{tripinfo}); resumen de estadísticas \\
Espera detenida & segundos con velocidad de a lo sumo 0.1~m/s (s/veh) & \texttt{waitingTime} \\
Paradas & veces que la velocidad baja de 0.1~m/s (por vehículo) & \texttt{waitingCount} \\
Cola promedio y máxima & detenidos sumados sobre las aproximaciones semaforizadas: media por
  segundo y máximo instantáneo (veh) & TraCI \\
\textit{Throughput} & vehículos que completan su recorrido dentro de la hora simulada &
  \texttt{tripinfo} \\
Emisiones & CO$_2$ (g/veh) y NO$_x$ (mg/veh) acumulados en el viaje & dispositivo de emisiones \\
No terminados & viajes sin llegada al cerrar la simulación & \texttt{tripinfo} \\
No insertados & vehículos que al cerrar la simulación seguían esperando entrar a la red (veh) &
  resumen de estadísticas \\
\midrule
Viaje, retraso y paradas de corredor & medias sobre los vehículos que recorren la avenida de
  extremo a extremo, en s/veh (las paradas, por vehículo); el retraso, también solo en el sentido
  de la punta & \texttt{tripinfo} \\
Progresión & fracción de esos vehículos que se detiene a lo sumo una vez & \texttt{tripinfo} \\
Cola máxima del enlace & máximo por carril de detenidos en los enlaces entre semáforos (veh) &
  TraCI \\
\textit{Spillback} & tiempo en que algún carril de un enlace entre semáforos tiene detenidos al
  menos un umbral de sus plazas & TraCI \\
\bottomrule
\end{tabular}
\caption[Métricas de desempeño]{Métricas de desempeño. El retraso total solo se usa en el
análisis de robustez. Las del bloque inferior se refieren a la avenida y a los enlaces entre
semáforos. El umbral de \textit{spillback} y el detalle de cada medición están en el
Capítulo~\ref{cap:metodologia}.}
\label{tab:metricas}
\end{table}
% fuente: capitulo2 y modulo U3 del capitulo4 originales (7bd8685); comun.py (leer_tripinfo: retraso_total_prom,
% retraso_insercion_prom, no_insertados; args_estadisticas)
% TODO(reestructura): citar la documentacion de las salidas tripinfo (timeLoss, departDelay, waitingTime,
% waitingCount) y statistic-output (waiting, departDelayWaiting) de SUMO.

El retraso es la métrica primaria e incluye frenadas, detenciones y aceleraciones. La espera
detenida es una parte del retraso, no su equivalente: un control puede reducirla y aumentar las
paradas sin mejorar el retraso. Retraso, colas y paradas son medidas habituales en la evaluación
de controles de tráfico \citep{FHWA_Toolbox_Definitions,FHWA_Timing_Manual_Chapter3}; las
emisiones las estima el modelo de emisiones de SUMO \citep{SUMO_Emissions}. El
\textit{throughput}, los no terminados y los no insertados solo discriminan entre controles cuando
la demanda supera la capacidad (Sección~\ref{sec:teoria-control}).

El retraso se cuenta desde que el vehículo entra a la red, y un vehículo entra solo si hay espacio
en el primer tramo de su ruta. Cuando la cola llega hasta el borde de la red, los que siguen
llegando esperan fuera, y esa espera no entra en el \texttt{timeLoss}
(Sección~\ref{sec:teoria-simulacion}). Es, sin embargo, retraso real: en una red de verdad esos
vehículos harían cola aguas arriba, en la vía que la alimenta, y cuánto esperan depende de cómo
el semáforo despacha la cola que tienen delante. Un control puede entonces bajar el
\texttt{timeLoss} dejando que la cola llegue al borde de la red, porque desde ahí la cola crece
fuera, donde el \texttt{timeLoss} no la ve. Para medir cuánto pesa ese efecto, el análisis de
robustez del Capítulo~\ref{cap:resultados} usa el retraso total, que suma las dos partes:
\begin{equation}
  D = \frac{\sum_{v \in V} \bigl(l_v + e_v\bigr) + W\, e_W}{|V| + W},
  \label{eq:retraso-total}
\end{equation}
donde $V$ es el conjunto de vehículos con registro en el \texttt{tripinfo}, incluidos los que no
terminaron su viaje al cerrar la simulación; $l_v$ es el retraso en la red del vehículo $v$
(\texttt{timeLoss}) y $e_v$, su espera para entrar (\texttt{departDelay}). $W$ es el número de
vehículos que al cerrar seguían esperando entrar, que no tienen registro y salen del resumen de
estadísticas, y $e_W$ es la espera media que llevaban acumulada en ese instante. Esos vehículos
seguirían esperando después del cierre, así que $W\,e_W$ es una cota inferior de su espera; lo
mismo vale para el retraso en la red de los que no terminaron, que se cuenta hasta el cierre. Si
ningún vehículo espera para entrar, $D$ coincide con el \texttt{timeLoss} medio. En el retraso
total la espera cuenta dentro y fuera de la red, y desplazar la cola no la hace desaparecer. La
espera para entrar se reporta también por separado, como métrica secundaria.

Conviene precisar la analogía con la ingeniería de transporte. El \textit{control delay} del HCM y
de la FHWA es el tiempo de viaje adicional atribuible al dispositivo de control
\citep{FHWA_Timing_Manual_Chapter3,HCM2000_or_2010}. El \texttt{timeLoss} de SUMO difiere de él
en dos sentidos. Incluye tiempo perdido que el semáforo no causa, como el de los giros y las
variaciones de la conducción; por eso un cambio porcentual medido sobre el \texttt{timeLoss}
subestima el efecto del control, ya que parte del retraso medido no depende de él. Y deja fuera la
espera para entrar a la red, que el semáforo sí puede causar; cuando la demanda supera la capacidad
de las entradas, el cambio medido sobre el \texttt{timeLoss} puede entonces sobrestimar ese
efecto. El retraso total corrige lo segundo, no lo primero.

Cada corrida, con cualquier control, registra además la recompensa de colas y tres retornos por
segundo, que se definen con la formulación (Capítulo~\ref{cap:metodologia}). Como se miden igual en
el fijo, el actuado y los aprendices, cada modelo puede compararse con los demás controles en su
propia recompensa. Son un diagnóstico y no entran en el veredicto.

\section{Aprendizaje por refuerzo}
\label{sec:rl}

\subsection{Proceso de decisión de Markov}

El aprendizaje por refuerzo (RL, \textit{reinforcement learning}) estudia a un agente que aprende
a decidir interactuando con un entorno. En cada paso $t$ observa el estado $s_t$ y elige una
acción $a_t$; el entorno responde con una recompensa $r_{t+1}$ y un nuevo estado $s_{t+1}$. El
agente busca maximizar la recompensa acumulada, no la inmediata \citep{SuttonBarto2018}.

La interacción se modela como un proceso de decisión de Markov (MDP), la tupla
$(\mathcal{S}, \mathcal{A}, P, R, \gamma)$: $\mathcal{S}$ es el conjunto de estados,
$\mathcal{A}$ el de acciones, $P(s' \mid s, a)$ la probabilidad de pasar al estado $s'$ al aplicar
$a$ en $s$, $R(s,a)$ la recompensa esperada de esa acción y $\gamma \in [0,1]$ el factor de
descuento (Ecuación~\ref{eq:retorno}). Se cumple la propiedad de Markov: el estado y la recompensa
siguientes dependen solo del estado y la acción actuales, no de la historia anterior
\citep{SuttonBarto2018}. Cuando no hace falta el índice temporal se escribe $s$, $a$, $r$ y $s'$
por $s_t$, $a_t$, $r_{t+1}$ y $s_{t+1}$.

En un semáforo, el estado resume el tráfico que ve el agente, la acción decide si se mantiene el
verde o se pasa a la fase siguiente y la recompensa penaliza la congestión; su forma concreta se
formula en el Capítulo~\ref{cap:metodologia}. El MDP supone que todos los pasos pesan lo mismo en
el descuento. Si las decisiones duran intervalos distintos, como cuando cambiar de fase consume
además el ámbar, el problema es en rigor semi-markoviano \citep{SuttonBarto2018}; el plan de pruebas trata
cada decisión como un paso con el mismo descuento (Capítulo~\ref{cap:implementacion}), y el análisis exploratorio posterior prueba descontar por segundo (Sección~\ref{sec:correccion-entrenamiento}).
% TODO(reestructura): verificar que SuttonBarto2018 (sec. 17.2) trate los procesos semi-markovianos; si no,
% citar Sutton, Precup y Singh (1999), que no esta en referencias.bib.

\subsection{Retorno, descuento y episodio}

La recompensa acumulada se mide con el retorno descontado
\begin{equation}
  G_t = r_{t+1} + \gamma\, r_{t+2} + \gamma^2 r_{t+3} + \cdots = \sum_{k=0}^{T-t-1} \gamma^k\, r_{t+k+1},
  \label{eq:retorno}
\end{equation}
donde $T$ es el último paso del episodio y $\gamma$ pondera las recompensas futuras: con
$\gamma = 0$ solo cuenta la inmediata, y cuanto más cerca de 1, más pesan las lejanas
\citep{SuttonBarto2018}. Con una recompensa constante $r$ y horizonte ilimitado el retorno es
$r/(1-\gamma)$, de modo que $1/(1-\gamma)$ se lee como el horizonte efectivo, en decisiones, que
el agente toma en cuenta; el Capítulo~\ref{cap:implementacion} usa esa lectura para interpretar
el descuento de las redes.

Un episodio es una corrida completa del entorno, desde el estado inicial hasta un estado terminal.
En esta tesis es una hora de demanda simulada sobre el mismo escenario, más el tiempo necesario
para que la red se vacíe; la red vacía es el estado terminal. Si la simulación se corta antes por
un tope de duración, el episodio queda truncado: el último estado no es terminal y su valor futuro
no es cero, distinción que respetan los objetivos de DQN y PPO de esta tesis
(Secciones~\ref{sec:dqn} y~\ref{sec:ppo}). El agente conserva lo aprendido de un episodio al siguiente.

\subsection{Política y funciones de valor}

Una política $\pi$ puede ser determinista, una acción por estado, o estocástica, una probabilidad
$\pi(a \mid s)$ de elegir $a$ en $s$. El valor de un estado es el retorno esperado desde él
siguiendo $\pi$, y el valor de acción, el retorno esperado si se aplica $a$ en $s$ y luego se
sigue $\pi$ \citep{SuttonBarto2018}:
\begin{equation}
  V^{\pi}(s) = \mathbb{E}_{\pi}\bigl[G_t \mid s_t = s\bigr], \qquad
  Q^{\pi}(s,a) = \mathbb{E}_{\pi}\bigl[G_t \mid s_t = s,\ a_t = a\bigr],
  \label{eq:valor}
\end{equation}
donde $\mathbb{E}_{\pi}$ es la esperanza cuando el agente sigue $\pi$ y $G_t$ es el retorno de la
Ecuación~\ref{eq:retorno}. El valor de acción óptimo $Q^*$, el de la política que maximiza el
valor en todos los estados, cumple la ecuación de optimalidad de Bellman
\begin{equation}
  Q^*(s,a) = R(s,a) + \gamma \sum_{s' \in \mathcal{S}} P(s' \mid s,a)\, \max_{a'} Q^*(s',a'),
  \label{eq:bellman}
\end{equation}
con $R$, $P$ y $\gamma$ los del MDP y el máximo sobre las acciones $a'$ del estado siguiente $s'$:
el valor de una acción es su recompensa esperada más el valor descontado de actuar de forma óptima
desde el estado siguiente \citep{SuttonBarto2018}. Conocida $Q^*$, la política óptima es voraz
(\textit{greedy}): elige en cada estado la acción de mayor valor, $\arg\max_a Q^*(s,a)$.
Q-learning y DQN estiman $Q^*$ a partir de la experiencia, sin conocer $P$ ni $R$.

\subsection{Exploración y política congelada}

Para estimar el valor de todas las acciones, el agente debe probarlas aunque parezcan peores que
la mejor conocida: es el compromiso entre explorar y explotar. La regla $\varepsilon$-greedy elige
con probabilidad $\varepsilon$ una acción al azar y, si no, la voraz, la de mayor valor estimado
\citep{SuttonBarto2018}. Mnih et al. reducen $\varepsilon$ linealmente durante el entrenamiento,
de modo que el agente explora al principio y explota al final \citep{Mnih2015_DQN}; el calendario
de cada aprendiz está en el Capítulo~\ref{cap:implementacion}. Los métodos con política
estocástica, como PPO, exploran muestreando la acción de $\pi(a \mid s)$, y su objetivo premia la
entropía de la política (Sección~\ref{sec:ppo}).

Mientras entrena, el agente explora y sus parámetros (la tabla de valores o los pesos de la red)
cambian dentro de cada episodio o entre episodios, así que las métricas de esos episodios
describen el aprendizaje, no el desempeño del controlador. Para medir este último se
\textit{congela} la política: se dejan de actualizar los parámetros y el agente deja de explorar.
Elige la acción de mayor valor en Q-learning, SARSA y DQN (equivale a $\varepsilon = 0$) y la de
mayor probabilidad en IPPO y MAPPO. En los tabulares, un empate entre las dos acciones se resuelve
al azar, con un generador sembrado para que la corrida se repita igual
(Capítulo~\ref{cap:implementacion}). La política congelada se corre igual que los controles de
referencia, y solo esos valores son comparables con el fijo y el actuado.

\section{Métodos tabulares: Q-learning y SARSA}
\label{sec:tabular}

\subsection{Diferencias temporales}

Los métodos tabulares guardan un valor $Q(s,a)$ por cada par de estado discreto y acción y
aprenden por diferencias temporales (TD): no esperan al final del episodio para conocer el
retorno, sino que lo aproximan con la recompensa inmediata más el valor estimado del estado
siguiente. La diferencia entre ese objetivo y la estimación actual es el error TD,
\begin{equation}
  \delta_t = r_{t+1} + \gamma\, V(s_{t+1}) - V(s_t),
  \label{eq:td}
\end{equation}
donde $V$ es la estimación vigente del valor de estado. La estimación se mueve una fracción de ese
error, $V(s_t) \leftarrow V(s_t) + \alpha\,\delta_t$, con $\alpha \in (0,1]$ la tasa de
aprendizaje \citep{SuttonBarto2018}. El mismo error reaparece en el crítico y en la ventaja de PPO
(Sección~\ref{sec:ppo}).

\subsection{Q-learning y SARSA}

Q-learning aplica esta idea al valor de acción, con el mejor valor del estado siguiente como
objetivo:
\begin{equation}
  Q(s,a) \leftarrow Q(s,a) + \alpha \bigl[r + \gamma\, \max_{a'} Q(s',a') - Q(s,a)\bigr],
  \label{eq:q-learning}
\end{equation}
donde $\alpha$ es la tasa de aprendizaje, $\gamma$ el factor de descuento, $r$ la recompensa
obtenida al ejecutar $a$ en $s$, $s'$ el estado siguiente y el máximo recorre las acciones $a'$ de
$s'$. El término $r + \gamma \max_{a'} Q(s',a')$ es el objetivo de diferencia temporal, una
estimación del retorno que se apoya en la propia tabla, no un retorno observado
\citep{SuttonBarto2018}. Q-learning es \textit{off-policy}: el objetivo usa la mejor acción de
$s'$ con independencia de la que se ejecute después, así que el agente puede explorar al azar y
aprender igualmente el valor de la política voraz. Si todos los pares de estado y acción se siguen
visitando y la tasa de aprendizaje decrece de forma adecuada, converge a $Q^*$ con probabilidad uno
\citep{SuttonBarto2018}. Son condiciones asintóticas que un entrenamiento finito no cumple; el
sentido operativo de convergencia de esta tesis está en la Sección~\ref{sec:evaluacion}.

SARSA es la versión \textit{on-policy}: elige primero la acción siguiente $a'$ con la misma
política $\varepsilon$-greedy con que actúa y actualiza con su valor,
\begin{equation}
  Q(s,a) \leftarrow Q(s,a) + \alpha \bigl[r + \gamma\, Q(s',a') - Q(s,a)\bigr],
  \label{eq:sarsa}
\end{equation}
donde $a'$ es la acción efectivamente seleccionada en $s'$, la que se ejecuta a continuación, y
$\alpha$, $\gamma$, $r$ y $s'$ son los de la Ecuación~\ref{eq:q-learning}. El nombre resume la
tupla $(s, a, r, s', a')$ de cada actualización \citep{SuttonBarto2018}. Al usar el valor de la
acción que la política realmente toma, incluida la exploración, SARSA mantiene la coherencia entre
exploración y actualización y suele dar comportamientos más conservadores que Q-learning.

\subsection{Discretización del estado}

Una tabla necesita un conjunto finito de estados, así que las magnitudes del tráfico, como el
número de vehículos detenidos, se agrupan en unos pocos niveles. Es una agregación de estados:
todos los estados de un grupo comparten un mismo valor \citep{SuttonBarto2018}. La discretización
hace viable la tabla y, a la vez, es su límite: el número de estados crece de forma multiplicativa
con las variables y sus niveles, y cuando el espacio es grande el agente visita solo una fracción
de él, de modo que la política elige entre valores apenas actualizados. Los niveles del estado
tabular de esta tesis están en el Capítulo~\ref{cap:metodologia}.

\section{Aproximación de funciones: DQN y Double DQN}
\label{sec:dqn}

\subsection{Redes neuronales como aproximadores}

Cuando el estado es continuo o demasiado grande para una tabla, el valor se representa con una
función parametrizada $Q(s,a;\theta)$, cuyo vector de parámetros $\theta$ tiene muchas menos
componentes que estados hay. Un cambio en $\theta$ modifica a la vez el valor de muchos estados,
lo que permite generalizar entre estados parecidos. Una red neuronal es una de esas funciones:
capas de combinaciones lineales seguidas de activaciones no lineales. Sus parámetros se ajustan
minimizando una pérdida por descenso de gradiente estocástico, un paso pequeño en contra del
gradiente de la pérdida en cada ejemplo \citep{SuttonBarto2018}. DQN y PPO promedian ese gradiente
sobre un minilote de ejemplos \citep{Mnih2015_DQN,Schulman2017_PPO}. Las redes de esta tesis
usan el optimizador Adam y recortan la norma del gradiente: si supera un umbral, el gradiente se
reescala hasta él (Capítulo~\ref{cap:implementacion}). Este recorte es distinto del recorte de la
razón de probabilidades de PPO (Sección~\ref{sec:ppo}).
% TODO(reestructura): citar el optimizador Adam (Kingma y Ba, 2015) y una fuente del recorte de la norma del
% gradiente; ninguna esta en referencias.bib.

\subsection{DQN}

DQN (\textit{Deep Q-Network}) aproxima $Q^*$ con una red neuronal entrenada para cumplir la
Ecuación~\ref{eq:bellman} sobre transiciones observadas \citep{Mnih2015_DQN}. Cada transición
$(s, a, r, s')$ se guarda en una memoria de repetición $\mathcal{D}$ de capacidad fija, y la red
se actualiza con minilotes muestreados al azar de ella, minimizando
\begin{equation}
  L(\theta) = \mathbb{E}_{(s,a,r,s') \sim \mathcal{D}}\Bigl[\bigl(y - Q(s,a;\theta)\bigr)^2\Bigr],
  \qquad y = r + \gamma \max_{a'} Q(s',a';\theta^-),
  \label{eq:dqn}
\end{equation}
donde $\theta$ son los parámetros de la red en línea, la que decide y se entrena, $y$ es el
objetivo y $\theta^-$ son los de la red objetivo, una copia de la red en línea que se actualiza
solo cada $C$ pasos de entrenamiento \citep{Mnih2015_DQN}. Si $s'$ es terminal, $y = r$; si el
episodio se truncó, $s'$ no es terminal y el objetivo conserva el término de la red objetivo. Si
hay acciones no permitidas en $s'$, el máximo recorre solo las permitidas (máscara de acciones,
Capítulo~\ref{cap:metodologia}).

La memoria de repetición rompe la correlación entre transiciones consecutivas y reutiliza cada
experiencia; la red objetivo mantiene fijo el objetivo durante $C$ pasos, de modo que no se
desplace con cada cambio de $\theta$. Ambos mecanismos estabilizan el aprendizaje. Mnih et al.
recortan además el error $y - Q(s,a;\theta)$ a $[-1, 1]$, lo que equivale a una pérdida
cuadrática para errores pequeños y lineal para los grandes \citep{Mnih2015_DQN}. Esa forma es la
pérdida de Huber que usa la implementación (Capítulo~\ref{cap:implementacion}).
% TODO(reestructura): citar la perdida de Huber (Huber, 1964) o la documentacion de smooth_l1_loss de PyTorch
% (rl_dqn.py usa F.smooth_l1_loss); no esta en referencias.bib.

\subsection{Double DQN}

El máximo de la Ecuación~\ref{eq:dqn} tiende a sobrestimar los valores: con estimaciones
ruidosas, usar el máximo de las estimaciones como estimación del máximo introduce un sesgo
positivo. Sutton y Barto llaman a este efecto sesgo de maximización y lo corrigen con el doble
aprendizaje, en el que una estimación elige la acción y otra, independiente, la valora
\citep{SuttonBarto2018}. Double DQN aplica esa idea con las dos redes de DQN: la red en línea
elige la acción del estado siguiente y la red objetivo la valora,
\begin{equation}
  y = r + \gamma\, Q\bigl(s',\, \arg\max_{a'} Q(s',a';\theta);\ \theta^-\bigr),
  \label{eq:double-dqn}
\end{equation}
con $r$, $\gamma$, $s'$, $\theta$ y $\theta^-$ como en la Ecuación~\ref{eq:dqn}. El DQN de esta
tesis es un Double DQN (Capítulo~\ref{cap:metodologia}).
% TODO(reestructura): citar Double DQN (van Hasselt, Guez y Silver, 2016); no esta en referencias.bib. La regla
% del objetivo esta en el Capitulo IV (parrafo de DQN) y en rl_dqn.py (actualizar).

\section{Métodos de gradiente de política: actor-crítico y PPO}
\label{sec:ppo}

\subsection{Gradiente de política y actor-crítico}

Los métodos de gradiente de política representan la política directamente, $\pi_\theta(a \mid s)$
con parámetros $\theta$, y ajustan $\theta$ en la dirección que aumenta el retorno esperado
\citep{SuttonBarto2018}. El teorema del gradiente de política permite estimar esa dirección con la
experiencia; en la forma que usan Schulman et al.,
\begin{equation}
  \hat g = \hat{\mathbb{E}}_t\bigl[\nabla_\theta \log \pi_\theta(a_t \mid s_t)\, \hat A_t\bigr],
  \label{eq:gradiente-politica}
\end{equation}
donde $\hat{\mathbb{E}}_t$ es el promedio sobre un lote de pasos de las trayectorias recogidas y
$\hat A_t$ estima la ventaja de la acción $a_t$ \citep{Schulman2017_PPO}. La ventaja
$A^\pi(s,a) = Q^\pi(s,a) - V^\pi(s)$ mide cuánto mejor es $a$ que el comportamiento medio de la
política en $s$; el gradiente sube la probabilidad de las acciones con ventaja positiva y baja la
de las otras.

Restar al retorno una línea base que solo depende del estado no cambia el valor esperado del
gradiente y puede reducir su varianza; la línea base natural es el valor del estado. En los
métodos actor-crítico una segunda función, el crítico $V_w(s)$ con parámetros $w$, estima ese
valor y se entrena con el error TD de la Ecuación~\ref{eq:td}; la política, llamada actor, se
actualiza con ventajas calculadas a partir de él \citep{SuttonBarto2018}. PPO es un método
actor-crítico: el actor decide y el crítico solo se usa para entrenar.

\subsection{Estimación de la ventaja}

PPO estima la ventaja con una versión truncada de la estimación generalizada de la ventaja (GAE,
\textit{generalized advantage estimation}), que suma los errores TD siguientes con pesos
decrecientes \citep{Schulman2017_PPO}:
\begin{equation}
  \hat A_t = \sum_{l=0}^{T-t-1} (\gamma\lambda)^l\, \delta_{t+l},
  \label{eq:gae}
\end{equation}
donde $\delta_{t+l}$ es el error TD de la Ecuación~\ref{eq:td} calculado con el crítico, $T$ el
último paso de la trayectoria, $\gamma$ el descuento y $\lambda \in [0,1]$ el parámetro que decide
cuánto pesan los pasos lejanos. Con $\lambda = 0$ la ventaja es el error TD de un paso, que se
apoya por completo en el crítico; con $\lambda = 1$ es el retorno observado hasta el final,
completado con el valor del último estado, menos el valor estimado de $s_t$
\citep{Schulman2017_PPO}. Los valores intermedios interpolan entre ambos, como el
$\lambda$-retorno de Sutton y Barto \citep{SuttonBarto2018}. Si la trayectoria acaba en un estado
terminal su valor es cero; si se trunca, el último término usa el valor que estima el crítico
(Capítulo~\ref{cap:implementacion}).
% TODO(reestructura): citar GAE (Schulman, Moritz, Levine, Jordan y Abbeel, 2016) para el compromiso entre sesgo
% y varianza que regula lambda; no esta en referencias.bib.

\subsection{Objetivo recortado de PPO}

PPO (\textit{Proximal Policy Optimization}) reutiliza cada lote de experiencia en varias pasadas de
optimización sin que la política se aleje demasiado de la que generó los datos
\citep{Schulman2017_PPO}. Para ello compara la política nueva con la anterior mediante la razón de
probabilidades
\begin{equation}
  \rho_t(\theta) = \frac{\pi_\theta(a_t \mid s_t)}{\pi_{\theta_{\text{old}}}(a_t \mid s_t)},
  \label{eq:ppo-razon}
\end{equation}
donde $\theta_{\text{old}}$ son los parámetros con que se recogió el lote. Schulman et al. la
llaman $r_t(\theta)$; aquí se escribe $\rho_t$, como la razón de muestreo por importancia de Sutton
y Barto, para no confundirla con la recompensa. El objetivo recortado es
\begin{equation}
  L^{\text{CLIP}}(\theta) = \hat{\mathbb{E}}_t\Bigl[\min\bigl(\rho_t(\theta)\, \hat A_t,\
  \operatorname{clip}\bigl(\rho_t(\theta),\, 1-\epsilon_{\text{clip}},\, 1+\epsilon_{\text{clip}}\bigr)\, \hat A_t\bigr)\Bigr],
  \label{eq:ppo-clip}
\end{equation}
donde $\hat A_t$ es la ventaja de la Ecuación~\ref{eq:gae}, $\operatorname{clip}$ limita la razón
al intervalo $[1-\epsilon_{\text{clip}}, 1+\epsilon_{\text{clip}}]$ y $\epsilon_{\text{clip}}$ es
el parámetro de recorte (el $\epsilon$ de Schulman et al., con subíndice para no confundirlo con la
$\varepsilon$ de exploración). El mínimo hace del objetivo una cota pesimista del objetivo sin
recortar: cuando la razón sale del intervalo en la dirección que lo mejoraría, el término deja de
crecer y su gradiente se anula, así que no hay incentivo para alejarse más de la política anterior
\citep{Schulman2017_PPO}.

Si actor y crítico comparten parámetros, el objetivo suma al término recortado el error del
crítico y un bono de entropía:
\begin{equation}
  L(\theta) = \hat{\mathbb{E}}_t\Bigl[L^{\text{CLIP}}_t(\theta)
  - c_1 \bigl(V(s_t) - V_t^{\text{targ}}\bigr)^2 + c_2\, S[\pi_\theta](s_t)\Bigr],
  \label{eq:ppo-total}
\end{equation}
donde $L^{\text{CLIP}}_t$ es el término de la Ecuación~\ref{eq:ppo-clip} en el paso $t$, $V(s_t)$
el valor que estima el crítico, $V_t^{\text{targ}}$ el retorno estimado que el crítico debe
aproximar, $c_1$ el coeficiente del valor, $S[\pi_\theta](s_t)$ la entropía de la distribución de
acciones en $s_t$ y $c_2$ el coeficiente de entropía; el objetivo se maximiza
\citep{Schulman2017_PPO}. El bono de entropía favorece distribuciones de acciones menos
concentradas y asegura una exploración suficiente. En esta tesis actor y crítico son redes
separadas, pero se optimizan juntas con esta suma y un solo optimizador; el error del crítico entra
como $\tfrac{1}{2}\bigl(V(s_t) - V_t^{\text{targ}}\bigr)^2$ (Capítulo~\ref{cap:implementacion}).

El algoritmo alterna dos etapas: recoge trayectorias con la política vigente
$\pi_{\theta_{\text{old}}}$ y calcula sus ventajas; luego optimiza el objetivo durante varias épocas
sobre minilotes de ese lote y hace $\theta_{\text{old}} \leftarrow \theta$
\citep{Schulman2017_PPO}. En esta tesis el lote es un episodio completo de todos los semáforos.
El tamaño de cada actualización se mide con la divergencia de Kullback--Leibler
$\mathrm{KL}\bigl[\pi_{\theta_{\text{old}}}(\cdot \mid s_t),\, \pi_\theta(\cdot \mid s_t)\bigr]$
entre la política anterior y la nueva, que Schulman et al. proponen también penalizar como
alternativa al recorte; aquí es solo un diagnóstico (Capítulo~\ref{cap:implementacion}).

\section{Aprendizaje por refuerzo multiagente}
\label{sec:marl}

\subsection{Observación parcial}

Con varios semáforos, varios agentes actúan a la vez sobre el mismo tráfico. Yu et al. formalizan
el caso cooperativo como un proceso de decisión de Markov descentralizado y parcialmente observable
(Dec-POMDP), $\langle \mathcal{S}, \mathcal{A}, O, R, P, n, \gamma \rangle$: $n$ agentes comparten
el espacio de acciones $\mathcal{A}$; en el estado global $s \in \mathcal{S}$, el agente $i$ solo
recibe una observación local $o_i = O(s; i)$; la transición $P(s' \mid s, \mathbf{a})$ depende de
la acción conjunta $\mathbf{a} = (a_1, \ldots, a_n)$; todos reciben la misma recompensa
$R(s, \mathbf{a})$, y $\gamma$ es el descuento. Cada agente decide con su propia observación,
$\pi_\theta(a_i \mid o_i)$, y entre todos maximizan el retorno descontado común
\citep{Yu2021_MAPPO}.

Lo que ve cada semáforo (sus colas, su fase y, en U4, los mensajes de sus vecinos) es entonces una
observación local $o_i$, no el estado de toda la red; también lo es la entrada del agente tabular,
que el Capítulo~\ref{cap:metodologia} llama estado con la terminología del MDP. Además, en esta tesis cada
agente recibe una recompensa local, calculada sobre sus propias aproximaciones y, en la variante
\emph{spill}, también sobre los enlaces por los que descarga hacia sus vecinos
(Ecuación~\ref{eq:recompensa-u4}), y no una recompensa común: el problema no es exactamente el
Dec-POMDP de recompensa compartida de Yu et al.
% TODO(reestructura): citar una fuente que formalice el caso con recompensas individuales por agente (juego
% estocastico parcialmente observable); no esta en referencias.bib.

\subsection{Agentes independientes y pesos compartidos}

En la configuración de agentes independientes, cada semáforo aprende con su observación y su
recompensa locales, sin intercambiar información, como si los demás fueran parte del entorno. Cada
agente percibe entonces un entorno no estacionario, porque los vecinos también cambian su
comportamiento mientras aprenden \citep{Yu2021_MAPPO}. Los agentes tabulares de esta tesis, con
una tabla por semáforo, son independientes en este sentido.

Si los agentes son homogéneos, pueden compartir parámetros: todos usan la misma red,
$\pi_\theta(a_i \mid o_i)$, y cada uno decide con su propia observación \citep{Yu2021_MAPPO}. Una
sola política aprende así de la experiencia de todos los semáforos. Para que la red distinga a los
agentes se le puede dar un identificador del agente, como hacen Yu et al. en el estado global de
su crítico \citep{Yu2021_MAPPO}. IPPO, DQN y MAPPO de esta tesis comparten parámetros entre los
semáforos, y su observación incluye el identificador del semáforo
(Capítulo~\ref{cap:metodologia}).

\subsection{IPPO, MAPPO y comunicación entre vecinos}

IPPO (\textit{Independent PPO}) es PPO con agentes independientes: el actor y el crítico de cada
agente solo reciben su observación local, $\pi_\theta(a_i \mid o_i)$ y $V(o_i)$, con parámetros
compartidos o no. Yu et al. lo usan como referencia de comparación \citep{Yu2021_MAPPO}.
% TODO(reestructura): citar IPPO (de Witt et al., 2020); no esta en referencias.bib.

MAPPO (\textit{Multi-Agent PPO}) conserva el actor descentralizado y cambia el crítico: durante el
entrenamiento estima el valor a partir de un estado global, $V(s)$, con más información que la
observación de cualquier agente, lo que busca dar estabilidad al aprendizaje en entornos
cooperativos. Una de las formas de estado global que discuten Yu et al. es la concatenación de las
observaciones locales, $(o_1, \ldots, o_n)$. Es el esquema de entrenamiento centralizado y
ejecución descentralizada: la información global solo se usa para entrenar, y en ejecución cada
agente decide con su actor y su observación, como en IPPO \citep{Yu2021_MAPPO}. El estado global
del MAPPO de esta tesis se describe en el Capítulo~\ref{cap:implementacion}.

Otra vía de coordinación es la comunicación explícita: cada agente recibe mensajes de sus vecinos
y los incorpora a su observación, de modo que la política sigue siendo descentralizada pero ve
información que su detector propio no puede medir. Es la vía de U4: en IPPO vecinos, IPPO spill,
DQN y MAPPO la coordinación entra por la observación y, en IPPO spill, que penaliza el
\textit{spillback}, también por la recompensa (Capítulo~\ref{cap:metodologia}).
% TODO(reestructura): citar una referencia de RL multiagente con comunicacion entre agentes; Yu2021_MAPPO no la
% trata y no hay otra en referencias.bib.

\section{Evaluación experimental}
\label{sec:evaluacion}

Comparar un controlador aprendido con los de referencia exige separar lo que se usa para aprender,
para elegir y para medir. Esta sección presenta los conceptos; el protocolo concreto, con sus
semillas y sus reglas, está en el Capítulo~\ref{cap:metodologia}.

\subsection{Unidad experimental, semillas y selección de la política}

Los episodios de entrenamiento no son observaciones independientes: el agente conserva lo aprendido
de uno al siguiente y su política cambia entre ellos. No pueden tratarse como réplicas
estadísticas. La unidad experimental es la corrida de evaluación, una simulación completa con la
política congelada (Sección~\ref{sec:rl}) sobre una semilla dada.

Las semillas se reparten en conjuntos disjuntos, cada uno con un uso: ajuste de los programas
fijos, entrenamiento, validación (para elegir la política) y evaluación (para medirla). Así la
política no se mide sobre las realizaciones con que aprendió ni con las que se eligió, y las
semillas de evaluación no intervienen en ninguna elección.

El desempeño durante el entrenamiento no tiene por qué mejorar de forma monótona, y la política del
último episodio no es necesariamente la mejor. Por eso se guardan los parámetros a intervalos
regulares, se valida cada copia con la política congelada y se elige la de mejor valor de
validación: una selección de modelo por validación, análoga a la parada temprana. Como el
resultado depende también de la semilla del propio entrenamiento, cada modelo se entrena varias
veces, cada una con su propia semilla base (Capítulo~\ref{cap:metodologia}), y el rango entre esos
entrenamientos mide su fragilidad. El
Capítulo~\ref{cap:metodologia} define la regla que elige la política que representa a cada modelo.
% TODO(reestructura): citar una fuente de seleccion de modelo por validacion o parada temprana y (opcional) otra
% de la variabilidad entre semillas de entrenamiento en RL profundo; no estan en referencias.bib.

La convergencia se juzga sobre la misma serie de validación. No es la convergencia asintótica de la
Sección~\ref{sec:tabular}, sino un criterio operativo: que la serie deje de tener tendencia y se
mantenga estable al final del entrenamiento, con los umbrales de la regla de veredicto
(Capítulo~\ref{cap:metodologia}).

\subsection{Diseño pareado, intervalos y pruebas}

Todos los controles se evalúan sobre las mismas semillas. Como la semilla fija el comportamiento de
los conductores (Sección~\ref{sec:teoria-simulacion}), dos controles corridos con la misma semilla
enfrentan el mismo tráfico, y la comparación puede hacerse por pares, semilla por semilla. El
diseño pareado reduce la variabilidad que se debe a la semilla y no al control.

Por control y métrica se reporta la media con su intervalo de confianza al 95\,\%, calculado con la
distribución $t$ de Student. Frente a una referencia, un modelo se compara por el cambio porcentual
de la media (Ecuación~\ref{eq:delta}), cuyo intervalo se estima por \textit{bootstrap} pareado
\citep{Efron1993_Bootstrap}: se remuestrean las semillas con reposición, conservando juntos los dos
valores de cada semilla, y se recalcula el cambio en cada remuestreo. La significancia se contrasta
con la prueba de Wilcoxon de rangos con signo \citep{Wilcoxon1945_Rangos} sobre las diferencias
pareadas por semilla, cuya hipótesis nula es que esas diferencias se reparten de forma simétrica
alrededor de cero.
% TODO(reestructura): citar el intervalo con t de Student; no esta en referencias.bib. Procedimiento:
% consolidar.py (docstring).

Las dos herramientas responden a preguntas distintas. El intervalo se refiere al cambio de la media,
que pueden mover unas pocas semillas con diferencias grandes; la prueba usa el signo y el rango de
todas las diferencias. Por eso pueden discrepar. Si unas pocas semillas tienen diferencias grandes
de un signo y las demás se reparten alrededor de cero, el intervalo puede excluir el cero sin que la
prueba sea significativa; si muchas diferencias pequeñas de un signo conviven con unas pocas
grandes del otro, puede ocurrir lo contrario. Lo primero ocurre en P5 y en P6 frente al fijo,
con IPPO vecinos en P5 y con IPPO spill en P6, y en P2 con MAPPO frente al actuado
(Capítulo~\ref{cap:resultados}). La regla de veredicto del Capítulo~\ref{cap:metodologia} exige
las dos condiciones.
% fuente: plan2/veredicto_congelado_0355.csv (= plan2/veredicto.csv del resultado principal; P5 ippo_vecinos:
% ic_fijo [0.4, 3.4], p_fijo 0.523; P2 mappo: ic_act [-4.5, -0.1], p_act 0.0687); P6: plan2/veredicto.csv
% (28.09.2026; ippo_spill: ic_fijo [0.6, 7.0], p_fijo 0.114; las demas celdas de P6 no discrepan)

El intervalo mide la variabilidad que introduce la semilla, y esa variabilidad depende de cómo se
declara la demanda. Si la demanda es determinista, con la salida programada de cada vehículo fijada de
antemano, la semilla cambia solo el comportamiento de los conductores y el intervalo no recoge la
variación de las llegadas. Es entonces condicional a la demanda declarada y subestima la
variabilidad de un tráfico real. La demanda de los escenarios de esta tesis es
determinista (Secciones~\ref{sec:escenarios} y~\ref{sec:estadistica}).

\subsection{Resumen sobre varios casos}

El veredicto se decide en cada caso por separado (Capítulo~\ref{cap:metodologia}). Para resumir un
modelo en todos los casos a la vez no sirve el promedio aritmético de sus retrasos, porque los casos
tienen escalas distintas y el de más retraso dominaría el promedio. Tampoco sirve el promedio
aritmético de los cocientes frente a una referencia, que puede ordenar dos sistemas de forma
distinta según la referencia elegida. La forma correcta de resumir resultados normalizados respecto
de una referencia es la media geométrica de los cocientes \citep{Fleming1986_MediaGeometrica}. No
depende de la escala de cada caso, porque cada cociente se toma dentro de su caso, y el cociente
entre las medias geométricas de dos modelos es el mismo con cualquier referencia. Es un resumen y no
reemplaza al veredicto por caso (Sección~\ref{sec:analisis-complementarios}).

Para saber de qué depende el resultado, si del caso o del modelo, la variación de la respuesta se
descompone con un análisis de varianza de dos factores con réplica. La suma de cuadrados total se
reparte entre el caso, el modelo, su interacción y la variación entre las réplicas de una misma
celda. La fracción que corresponde a cada fuente,
$\eta^2 = \mathrm{SC}_{\text{fuente}} / \mathrm{SC}_{\text{total}}$, con SC la suma de cuadrados,
es el eta cuadrado, una medida del tamaño del efecto \citep{Olejnik2003_Eta}. Un $\eta^2$ alto del
caso indica que el resultado depende sobre todo del escenario, y uno alto del modelo, que depende
sobre todo del algoritmo. Las fracciones suman uno y se leen como el reparto de la variación en esa
muestra de casos y modelos; su valor depende del diseño, así que no se compara con el de estudios
de otro diseño. La respuesta y las réplicas que usa esta tesis están en la
Sección~\ref{sec:analisis-complementarios}.
